\chapter{최적화를 깨지 않게 검증하는 방법}

\section{검증 대상은 값, cycle, protocol 세 층이다}

수학 결과만 같다고 hardware optimization이 끝난 것은 아니다. pipeline을 바꾼
뒤에는 다음 세 층을 따로 확인한다.

\begin{longtable}{L{28mm}L{61mm}L{72mm}}
\toprule
층 & 확인할 계약 & 이 사례의 예 \\
\midrule
Value & modular result와 transform coefficient & multiplier/butterfly/full NTT reference 비교 \\
Cycle & input, result, valid, last, address tag 정렬 & issue-to-writeback latency, layer done lookahead \\
Protocol & ready/valid, backpressure, reset, phase exclusivity & AXI gap/TREADY stall/TLAST, LOAD--CORE--DUMP \\
\bottomrule
\end{longtable}

\section{latency를 하나의 상수로 관리한다}

최종 core는 \rtlfile{XXX\_ISSUE\_TO\_WB\_LAT=13}을 definitions file에 한 번만
정의한다. top과 core가 서로 다른 숫자를 localparam으로 가지면 arithmetic
latency 수정 후 metadata pipeline이 조용히 어긋날 수 있다.

testbench에서는 다음 사건의 cycle 차이를 검사한다.

\begin{enumerate}
  \item address generator의 issue valid
  \item BRAM read request와 response
  \item butterfly input acceptance
  \item writeback valid/address/data
  \item final writeback과 layer/core done
\end{enumerate}

valid만 맞추면 충분하지 않다. A/B/C/D address tag와 두 zeta, radix-3의 D-lane
absence, inverse mode의 twiddle 선택이 같은 transaction을 가리켜야 한다.

\section{assertion으로 optimization 전제를 코드에 남긴다}

최종 RTL의 simulation-only check는 설계자가 area/timing을 줄일 때 사용한 전제를
검사한다.

\begin{itemize}
  \item transform 실행 중 crypto, $N$, inverse parameter가 바뀌지 않는다.
  \item physical row가 576-entry bank 범위를 넘지 않는다.
  \item 같은 BRAM row에 두 port가 동시에 write하지 않는다.
  \item multiplier operand와 intermediate signed range가 증명 범위 안에 있다.
  \item Montgomery의 $T-uq$ low 16 bit가 0이고 최종 값이 $[0,q)$다.
  \item dump FIFO count와 credit이 depth를 넘지 않는다.
\end{itemize}

\begin{lstlisting}
`ifndef SYNTHESIS
always @(posedge iSYS_CLK) begin
  if (response && (count_q == FIFO_DEPTH) && !pop)
    $error("unreserved BRAM response");
  if (credit_q > FIFO_DEPTH)
    $error("dump FIFO credit overflow");
end
`endif
\end{lstlisting}

assertion은 test vector가 우연히 문제 조건을 만들지 않아도 내부 계약 위반을
찾는다. \rtlfile{SYNTHESIS} guard로 production hardware에는 들어가지 않는다.

\section{변경 종류별 필요한 검증}

\begin{longtable}{L{44mm}L{118mm}}
\toprule
변경 & 최소 검증 \\
\midrule
Register 위치/latency 변경 & cycle scoreboard, valid/last/address alignment, back-to-back transaction \\
Bit width/signed 변경 & 경계값, random arithmetic, range assertion \\
Memory mapping 변경 & 전 주소 bijection, row bound, bank/port conflict, roundtrip \\
BRAM coding template 변경 & simulation read-during-write semantics와 synthesis inference report \\
FIFO/credit 변경 & random source gap와 destination backpressure, overflow/duplicate/loss check \\
Mode sharing 변경 & mode별 unit test와 mode-stability assertion \\
Scheduler 변경 & issue sequence, per-layer count, all transform roundtrip/reference \\
\bottomrule
\end{longtable}

\section{실행 가능한 regression}

Companion AXI4-Stream 저장소 root에서 최종 RTL regression을 실행한다.

\begin{lstlisting}[language=bash,numbers=none]
./tool_sim/run_all.sh
\end{lstlisting}

성공 시 다음 다섯 test가 각각 PASS로 표시된다.

\begin{enumerate}
  \item memory mapping: bijection과 radix-2 bank balance
  \item unified modular multiplier: 30 arithmetic cases
  \item unified butterfly: 14 radix-2/radix-3/trinomial cases
  \item full NTT/INTT: 지원하는 8 configuration
  \item AXI integration: source gap, backpressure, TLAST
\end{enumerate}

마지막 줄은 \texttt{ALL\_AXI\_STREAM\_TESTBENCHES\_PASS}다. 보존한 100~MHz
baseline은 다음처럼 실행한다.

\begin{lstlisting}[language=bash,numbers=none]
./old/tool_sim/run_all.sh
\end{lstlisting}

이 regression은 unified butterfly, ML-KEM radix-2 roundtrip, NTRU+768의
trinomial/radix-3/radix-2 roundtrip을 검사하고 3/3 PASS를 출력한다. 두 suite는
변경 전후 모두 실행 가능하므로 ``새 코드만 맞는다''가 아니라 보존 baseline의
동작도 재확인할 수 있다.

\section{simulation PASS와 implementation PASS를 구분한다}

Icarus regression은 기능과 protocol을 확인하지만 DSP48E1 placement, BRAM
inference, timing closure를 증명하지 않는다. 최종 acceptance에는 다음이 모두
필요하다.

\begin{itemize}
  \item RTL regression PASS
  \item synthesis에서 DSP 4개와 의도한 BRAM 확인
  \item post-route setup/hold failing endpoint 0
  \item bitstream 생성과 DRC error 0
  \item board KAT/reference 비교 PASS
  \item benchmark의 clock와 measurement interval 기록
\end{itemize}

이 사례의 packed-BRAM board binary는 benchmark 전에 6개 scheme의 NTT와 INTT,
총 12개 transform을 software reference와 비교했다. completed state가 기록된 뒤
성능을 측정했으므로 기능 검증과 timing measurement의 실행 순서가 분명하다.
